\documentclass[11pt]{article}
\usepackage[margin=0.9in]{geometry}
\usepackage{amsmath,amssymb,amsthm,booktabs,microtype}
\usepackage[colorlinks=true,allcolors=blue]{hyperref}
\newtheorem{proposition}{Proposition}
\newcommand{\E}{\mathbb E}
\newcommand{\PP}{\mathbb P}
\newcommand{\ind}{\mathbf 1}
\newcommand{\Lgen}{\mathcal L}
\DeclareMathOperator{\arcosh}{arcosh}
\title{Research Prizes with Uncertain Compute Costs:\\From Entry Games to Brownian Stopping Races}
\author{}
\date{September 2026}
\begin{document}
\maketitle
\begin{abstract}
We study competition for a research prize when the required computation is
uncertain and researchers have different costs. A one-period model gives
explicit mixed-entry conditions and an exact potential, hence a pure
equilibrium also exists. A continuous-time extension replaces committed
attempts by privately observed diffusions with irreversible abandonment;
the first solution wins. Against fixed rival strategies, their solution
hazard acts as additional discounting. We solve the constant-hazard,
driftless case in closed form and formulate the general equilibrium as
coupled stopping and forward-distribution equations. This connects research
races to Brownian bandit pruning, while distinguishing a strategic fixed
point from a planner's optimization problem.
\end{abstract}

\section{Setting and relation to the literature}
A large mathematical problem can be worth attempting alone even when many
potential solvers prefer to stay out of a prize contest. The relevant cost
is not necessarily a known token budget: it may be the random expenditure
incurred before success, abandonment, or a rival's solution. The question
is how competition changes participation and continued computation.

The one-period benchmark belongs to the literature on endogenous contest
entry. Morgan, Orzen, and Sefton~\cite{morgan} study entry and subsequent
rent-seeking, including the role of the outside option. Their post-entry
effort game differs from our fixed-attempt lottery. Fu, Jiao, and
Lu~\cite{fu} analyze complete-information lottery contests with costly
entry and subsequent bidding; we instead hold each attempt's success
probability fixed and allow researcher heterogeneity. Fullerton and
McAfee~\cite{fullerton} study heterogeneous research contestants and
selection through auctioned entry, rather than the uncoordinated binary
entry considered here. Our finite-game potential is in the tradition of
Rosenthal~\cite{rosenthal} and Monderer and Shapley~\cite{monderer}.
These connections locate the benchmark; we make no novelty claim for
entry indifference or the potential-game method.

The dynamic model belongs to the research-race tradition represented by
Loury~\cite{loury}. Its computational connection is to the Brownian bandit
framework~\cite{cotton}: a backward continuation problem and a forward
population calculation are coupled. Here the coupling is strategic and
the payoff is first passage to a solution, rather than a terminal maximum.
We prove the elementary entry and constant-hazard results below. The
general dynamic equilibrium is a characterization, not an existence,
uniqueness, or numerical convergence theorem.

\section{One-period entry with heterogeneous costs}
There are $n\ge2$ risk-neutral researchers and a prize $V>0$. Researcher
$i$ enters with probability $q_i$, incurs a random attempt cost with finite
mean $k_i>0$, and succeeds with probability $p_i\in(0,1]$. Parameters are
common knowledge. Entry randomizations and successes are independent
across researchers. Own cost and own success may be dependent. An entrant
commits to the entire attempt; nonentry pays zero. If several entrants
succeed, one is selected uniformly to receive $V$; if none succeeds, no
prize is paid. Assume $k_i<Vp_i$, so each attempt is profitable alone.

\begin{proposition}[Entry equilibrium]\label{prop:entry}
Define
\begin{equation}\label{eq:entry}
D_i(q_{-i})=Vp_i\int_0^1\prod_{j\ne i}
 (1-q_jp_j+q_jp_jz)\,dz-k_i.
\end{equation}
A profile $q\in[0,1]^n$ is a Nash equilibrium exactly when
\begin{equation}\label{eq:comp}
q_i=0\Rightarrow D_i\le0,\qquad
0<q_i<1\Rightarrow D_i=0,\qquad
q_i=1\Rightarrow D_i\ge0
\end{equation}
for every researcher.
\end{proposition}
\begin{proof}
Let $K_{-i}$ count successful opponents. Conditional on own success, the
prize probability is $(1+K_{-i})^{-1}$. Since
$1/(k+1)=\int_0^1z^k\,dz$, independence gives the integral in
\eqref{eq:entry}. Expected payoff is $q_iD_i(q_{-i})$; optimizing this
linear function on $[0,1]$ yields \eqref{eq:comp}.
\end{proof}

\paragraph{Two researchers.}
The fully mixed equilibrium, when it exists, is
\begin{equation}\label{eq:two}
q_1=\frac{2}{p_1}\left(1-\frac{k_2}{Vp_2}\right),\qquad
q_2=\frac{2}{p_2}\left(1-\frac{k_1}{Vp_1}\right).
\end{equation}
Both probabilities are interior exactly when
$Vp_i(1-p_j/2)<k_i<Vp_i$ for $i\ne j$. For $V=100$,
$p_1=p_2=1$, and $(k_1,k_2)=(60,80)$, the solution is
$(q_1,q_2)=(0.4,0.8)$. The cheaper researcher enters less often in this
mixed equilibrium because each probability makes the \emph{opponent}
indifferent. Both $(1,0)$ and $(0,1)$ are also equilibria. The example
does not imply a general negative relationship between efficiency and
participation across selected equilibria.

\paragraph{Symmetric benchmark.}
For common $p,k$, the gross entry return against common probability $q>0$ is
\begin{equation}\label{eq:symmetric}
R(q)=\frac{V[1-(1-pq)^n]}{nq},\qquad R(0)=Vp.
\end{equation}
It is strictly decreasing. Under $0<k<Vp$, the unique symmetric equilibrium
has $q=1$ if $k\le V[1-(1-p)^n]/n$; otherwise its interior probability
solves $R(q)=k$. For $p=1$, universal entry requires $k\le V/n$.
In an interior symmetric equilibrium,
$nqk=V[1-(1-pq)^n]$: expected expenditure equals expected prize payout.

\begin{proposition}[An exact potential]\label{prop:potential}
For a pure entrant set $A$, let $B_i$ be independent Bernoulli variables
with parameters $p_i$, and let $H_m=\sum_{\ell=1}^m1/\ell$, with $H_0=0$.
Then
\begin{equation}\label{eq:potential}
\Phi(A)=V\E\!\left[H_{\sum_{i\in A}B_i}\right]
       -\sum_{i\in A}k_i
\end{equation}
is an exact potential. In particular, a pure Nash equilibrium exists.
\end{proposition}
\begin{proof}
Adding $i\notin A$ changes \eqref{eq:potential} by
$Vp_i\E[(1+\sum_{j\in A}B_j)^{-1}]-k_i$, exactly the entrant's payoff
change. A maximizing entrant set exists because the set of profiles is
finite, and admits no profitable unilateral deviation.
\end{proof}
Thus mixing is not necessary for existence. When $p_i=1$, the potential
reduces to $VH_{|A|}-\sum_{i\in A}k_i$, the harmonic form associated with
a congested common resource. The potential is not social welfare:
aggregate expected payoff is
$V[1-\prod_{i\in A}(1-p_i)]-\sum_{i\in A}k_i$.

\section{A continuous-time race with private progress}
We now change the allocation rule: the \emph{first} verified solution wins.
Researcher $i$ starts at $x_i<b$ and, while working, observes
\begin{equation}\label{eq:diffusion}
dX_i(t)=\mu_i\,dt+\sigma_i\,dW_i(t),\qquad \sigma_i>0.
\end{equation}
The Brownian motions are independent. Hitting $b$ produces a solution;
working costs $c_i>0$ per unit time. Each researcher can abandon
irreversibly, but cannot pause, restart, or change compute intensity in
this baseline. Everyone discounts at rate $r\ge0$ and observes the first
solution announcement immediately. All parameters and initial-state
distributions are common knowledge; realized rival states and abandonment
decisions are unobserved. Initial states and private randomizations are
independent across researchers.

The state $X_i$ is a stylized measure of an approach's promise, not a claim
that proof completion has a literal observable percentage. The model
requires such a state to summarize the researcher's predictive information.
Common uncertainty about problem difficulty would violate independence
and require additional beliefs.

For each strategy, define its \emph{potential} solution time $T_i$ by
running it as though no rival ever announced success; set $T_i=\infty$
if it abandons first. Let $\theta_i$ be the corresponding abandonment
time and $T_{-i}=\min_{j\ne i}T_j$. Actual payoff is
\begin{equation}\label{eq:payoff}
\E\left[V e^{-rT_i}\ind_{\{T_i<T_{-i},\ T_i<\infty\}}
 -\int_0^{\theta_i\wedge T_i\wedge T_{-i}}e^{-rt}c_i\,dt\right].
\end{equation}
Nondegenerate diffusions started below $b$ have no finite-time ties in
this setting. Infinite times in \eqref{eq:payoff} contribute no prize.
Unlike Section~2, costs cease when a rival wins.

\subsection{Best response to a rival solution hazard}
Fix rival strategies and let
\[
h_{-i}(t)=-\frac{d}{dt}\log\PP(T_{-i}>t)
\]
where the survival probability is positive and differentiable. Independence
makes this a deterministic function of time, conditional on the prize
remaining available. Define
$\Lgen_i=\mu_i\partial_x+(\sigma_i^2/2)\partial_{xx}$ and
$D_i(t,s)=\exp\{-\int_t^s[r+h_{-i}(u)]\,du\}$.
Conditional on time $t$, own state $x$, and no prior solution, the value is
\begin{equation}\label{eq:stopvalue}
v_i(t,x)=\sup_{\theta\ge t}\E_{t,x}\left[
 VD_i(t,\tau_b)\ind_{\{\tau_b\le\theta,\ \tau_b<\infty\}}
 -\int_t^{\theta\wedge\tau_b}D_i(t,s)c_i\,ds\right],
\end{equation}
where $\tau_b$ is own first passage to $b$ and immediate abandonment is
allowed. Success takes precedence at $b$.

Dynamic programming gives the obstacle equation
\begin{equation}\label{eq:hjb}
\boxed{\max\{-v_i,\ \partial_t v_i+\Lgen_i v_i
 -(r+h_{-i})v_i-c_i\}=0,\qquad v_i(t,b)=V.}
\end{equation}
The value is the bounded optimal-stopping value in \eqref{eq:stopvalue};
where classical regularity fails, \eqref{eq:hjb} is interpreted in the
usual weak or viscosity sense. For a deadline $T$, add
$v_i(T,x)=0$ for $x<b$. Competition acts as additional discounting, but its
rate must ultimately be consistent with rival behavior.

\section{An exact constant-hazard benchmark}
Take one researcher with $\mu=0$, cost rate $c$, volatility $\sigma$,
an infinite horizon, and constant rival hazard $h$. Write $\rho=r+h$.

\begin{proposition}[Optimal abandonment boundary]\label{prop:boundary}
For $\rho>0$, define
\begin{equation}\label{eq:boundary}
\kappa=\frac{\sqrt{2\rho}}{\sigma},\qquad
a=b-\frac{1}{\kappa}\arcosh\left(1+\frac{\rho V}{c}\right).
\end{equation}
The optimal policy abandons at $x\le a$ and continues on $(a,b)$. Its value is
\begin{equation}\label{eq:closed}
v(x)=\begin{cases}
0,&x\le a,\\[2pt]
\displaystyle\frac{c}{\rho}\bigl[\cosh(\kappa(x-a))-1\bigr],&a<x<b,\\[2pt]
V,&x=b.
\end{cases}
\end{equation}
For $\rho=0$, the limiting solution has
$a=b-\sigma\sqrt{V/c}$ and $v(x)=c(x-a)^2/\sigma^2$ on $(a,b)$.
\end{proposition}
\begin{proof}
On $(a,b)$, solve $(\sigma^2/2)v''-\rho v-c=0$ with
$v(a)=v'(a)=0$. This yields \eqref{eq:closed}; $v(b)=V$ gives
\eqref{eq:boundary}. The function is nonnegative, bounded by $V$, and
continuously differentiable at $a$. In the stopping region the continuation
operator equals $-c<0$, while in the continuation region it equals zero.
Discounted It\^o verification therefore bounds every admissible stopping
payoff above by $v$, with equality for exit from $(a,b)$. For $\rho>0$ the
bounded terminal remainder vanishes; for $\rho=0$ the proposed interval
exit has finite expected duration and the quadratic candidate verifies
directly. Finally, $\arcosh(1+z)\sim\sqrt{2z}$ gives the stated limit.
\end{proof}

The continuation distance $b-a$ increases with $V$ and decreases with $c$
and $\rho$. For the last claim, put $z=\rho V/c$ and observe that
$\arcosh(1+z)/\sqrt z$ decreases for $z>0$. Indeed, writing
$y=\arcosh(1+z)$ reduces its derivative sign to
$2\tanh(y/2)<y$. Thus greater competition raises the abandonment threshold.
These are best-response comparative statics at a prescribed hazard, not
comparative statics of the full equilibrium. A constant positive hazard
is a benchmark; it need not be generated by a finite population of
abandoning researchers.

\section{Closing the equilibrium loop}
For deterministic time--state stopping policies, potential solution times
remain independent. Let
\begin{equation}\label{eq:survival}
F_i(t)=\PP(T_i\le t),\qquad f_i(t)=F_i'(t),\qquad S_i(t)=1-F_i(t).
\end{equation}
The distributions can be defective: $S_i(\infty)>0$ when some paths never
solve. Then
\begin{equation}\label{eq:hazard}
\boxed{h_{-i}(t)=\sum_{j\ne i}\frac{f_j(t)}{S_j(t)}.}
\end{equation}
To see this, $\PP(T_{-i}>t)=\prod_{j\ne i}S_j(t)$; logarithmic
differentiation proves \eqref{eq:hazard}. The same construction allows
independent private mixtures over policies.

\paragraph{Forward calculation.}
For a regular continuation region $(a_i(t),b)$, let $m_i(t,x)$ be the
unconditional density of the potential process still working and unsolved.
It satisfies
\begin{equation}\label{eq:forward}
\partial_t m_i=-\mu_i\partial_xm_i
 +\frac{\sigma_i^2}{2}\partial_{xx}m_i,
\qquad m_i(t,a_i(t))=m_i(t,b)=0,
\end{equation}
with the specified initial law and absorption at either boundary.
The solution density is the upper-boundary flux
\begin{equation}\label{eq:flux}
f_i(t)=-\frac{\sigma_i^2}{2}\partial_xm_i(t,b^-).
\end{equation}
For nonsmooth or jumping stopping regions, use the killed transition law
and account explicitly for newly abandoned mass instead of assuming
classical boundary regularity.

Two bookkeeping points matter. First, $S_i(t)$ includes abandoned paths;
it is \emph{not} $\int m_i(t,x)\,dx$. Second, \eqref{eq:forward} does not
kill mass at the rival hazard: it describes potential paths without rival
announcements. Including such killing here and also in \eqref{eq:hjb}
would count the rival's termination twice.

\paragraph{Equilibrium characterization.}
A consistent profile consists of strategies optimal for
\eqref{eq:stopvalue}, together with their forward laws and hazards satisfying
\eqref{eq:survival}--\eqref{eq:hazard}. Such a profile is a Nash equilibrium
under the stated private-information structure: fixing rivals fixes their
hazard, and each researcher already solves the resulting best-response
problem. Conversely, an equilibrium represented by these Markov stopping
policies must satisfy this consistency condition. In the symmetric case,
\begin{equation}\label{eq:symmetricdynamic}
h_{-i}(t)=(n-1)\frac{f(t)}{1-F(t)},
\end{equation}
so one backward problem and one forward law suffice. A finite number of
researcher types requires one pair per type, with the appropriate count of
opponents of each type. No large-population limit is needed for this reduction.

The baseline has no setup charge: immediate abandonment supplies the entry
option. With an additional setup cost $K_i$, entry compares the initial
continuation value with $K_i$. Mixed entry then requires equality, and
nonentry contributes mass at $T_i=\infty$. More generally, randomness of
Brownian outcomes is not strategic mixing; a deterministic stopping rule
can be optimal. Mixtures must be supported on best responses.

\section{Brownian bandits inside a strategic race}
The Brownian bandit model~\cite{cotton} uses independent, irreversibly
prunable paths, an expected path-time budget, and a terminal-maximum
objective. Its population-blind Poisson relaxation couples a backward
obstacle problem to a forward survivor calculation through terminal
pivotal value. The present game reuses this architecture, replacing the
planner's marginal terminal reward by a first-solution prize and replacing
the optimization coupling by rival-hazard consistency. Its concavity and
optimization certificates do not thereby establish equilibrium convergence.

One researcher can also operate a portfolio of approaches $X_{i\ell}$.
Let $A_i(t)$ be the active set, $C_i(A_i)$ its cost rate, and $\tau_i$ the
first solution among its active approaches. Against rival hazard $h_{-i}$,
the researcher optimizes
\begin{equation}\label{eq:portfolio}
\sup_{A_i}\E\left[
VD_i(0,\tau_i)\ind_{\{\tau_i<\infty\}}
-\int_0^{\tau_i}D_i(0,t)C_i(A_i(t))\,dt\right],
\end{equation}
where admissible controls include irreversible pruning and cost becomes
zero after total abandonment. A hard compute budget requires a remaining-
budget state. This nests an internal allocation problem inside an external
prize race. Correlations among approaches may be modeled internally, but
the simple product survival formula still requires independence across
researchers.

If compute intensity $u\in[0,\bar u]$ time-changes a single approach,
$dX=\mu u\,dt+\sigma\sqrt u\,dW$ and cost is $cu\,dt$, the continuation
operator becomes
\begin{equation}\label{eq:intensity}
\partial_t v-(r+h_{-i})v
+\sup_{0\le u\le\bar u}u(\Lgen v-c).
\end{equation}
Here $u=0$ permits pausing, so this is a separate extension of the baseline.
Linearity makes the best intensity an endpoint except at indifference.
A finite capacity or increasing marginal compute cost is necessary to avoid
an ill-posed instantaneous-acceleration incentive when $\Lgen v-c>0$.

\section{A numerical route and its limits}
A finite-deadline approximation suggests the following procedure: initialize
rival hazards; solve each backward stopping problem; propagate the induced
potential laws forward; recompute hazards using \eqref{eq:hazard}; and
iterate, with damping if useful. Fixed-point iteration is a proposed solver,
not a convergence result. In particular, this is not automatically the
concave optimization solved by the Brownian bandit implementation.

Validation should compare the stopping solver with
Proposition~\ref{prop:boundary}, check conservation of active, abandoned,
and solved mass, and measure each researcher's payoff gain from recomputing
a best response against the candidate rivals. That unilateral deviation
gain, rather than successive-iterate agreement alone, measures approximate
Nash behavior. Time, state, and horizon refinement remain necessary.

The main unresolved mathematical tasks are existence and selection of the
coupled equilibrium and conditions for convergence of its numerical scheme.
Publicly observed rival progress requires a joint-state game; correlated
difficulty requires filtering. Neither extension is captured by a
deterministic calendar-time hazard alone. The tractable model isolates the
central mechanism: the option to keep computing is exercised against an
endogenous risk that a competitor makes the prize unavailable.

\begin{thebibliography}{9}
\bibitem{fu}
Q. Fu, Q. Jiao, and J. Lu (2015).
Contests with endogenous entry. \emph{International Journal of Game Theory}
44(2), 387--424.
\href{https://doi.org/10.1007/s00182-014-0435-9}{doi:10.1007/s00182-014-0435-9}.

\bibitem{morgan}
J. Morgan, H. Orzen, and M. Sefton (2012).
Endogenous entry in contests. \emph{Economic Theory} 51, 435--463.
\href{https://doi.org/10.1007/s00199-010-0544-z}{doi:10.1007/s00199-010-0544-z}.

\bibitem{fullerton}
R. L. Fullerton and R. P. McAfee (1999).
Auctioning entry into tournaments. \emph{Journal of Political Economy}
107(3), 573--605.
\href{https://doi.org/10.1086/250072}{doi:10.1086/250072}.

\bibitem{rosenthal}
R. W. Rosenthal (1973).
A class of games possessing pure-strategy Nash equilibria.
\emph{International Journal of Game Theory} 2, 65--67.
\href{https://doi.org/10.1007/BF01737559}{doi:10.1007/BF01737559}.

\bibitem{monderer}
D. Monderer and L. S. Shapley (1996).
Potential games. \emph{Games and Economic Behavior} 14(1), 124--143.
\href{https://doi.org/10.1006/game.1996.0044}{doi:10.1006/game.1996.0044}.

\bibitem{loury}
G. C. Loury (1979).
Market structure and innovation. \emph{The Quarterly Journal of Economics}
93(3), 395--410.
\href{https://doi.org/10.2307/1883165}{doi:10.2307/1883165}.

\bibitem{cotton}
P. Cotton (2026).
Budgeted Brownian races: wavefront pruning under a path-time budget.
Project documentation, accessed September 9, 2026.
\url{https://brownianbandit.microprediction.org/}.
\end{thebibliography}
\end{document}
